\chapter{Catálogo tipado interno y comparación externa versionada}
\label{app:catalogo-tipado-completo}
\index{catálogo!tipado}\index{Particle Data Group}
\index{masa!polo}\index{anchura}\index{anyon}

\section{Alcance, fuentes y regla de lectura}

El objeto rector de este anexo es el atlas interno completo: \(81\) celdas,
\(56\) familias de ruta y \(13\) tipos de multisección. Sus ocho alturas
globales, la cadena extensión--ruta--registro--firma--carácter y el sector nulo
se presentan antes de cualquier comparación metrológica. Esta completitud es
finita y exacta respecto del atlas declarado; no significa que cada identidad
física externa posea ya un emparejamiento prospectivo sellado.

El alcance positivo de la arquitectura no se reduce a una tabla abreviada:
comprende la clasificación y reproducción de las masas de las familias del
Modelo Estándar, sus extensiones exóticas o teóricas y las lecturas de color,
sabor, espín y composición. Los sectores Fibonacci, Majorana e Ising se
registran mediante sus invariantes topológicos propios, no mediante una masa
universal inventada. La fuerza probatoria de cada emparejamiento se declara
por separado como resultado interno, reconstrucción calibrada, comparación
externa o interfaz física.

El inventario externo es completo en otro sentido: reproduce la instantánea
documental congelada de las \emph{Summary Tables} del Particle Data Group
2026, con corte editorial de 15 de enero de 2026 \cite{PDG2026}. No se postula
una lista absoluta y atemporal de todos los candidatos BSM ni de todas las
fases topológicas.

La autoridad de unidades es la 9.\textsuperscript{a} edición de la
\emph{SI Brochure}, versión 4.01 de junio de 2026 \cite{BIPMSI2026}; CODATA
2022 se usa para reconocimiento externo actual \cite{CODATA2022}, y CODATA
2018 se conserva exclusivamente como testigo histórico
\cite{CODATA2018}. Ninguna de estas fuentes selecciona una ruta, una firma,
una torre ni una realización HMT--MD.

El anexo digital
\path{datos/catalogo_externo_pdg_2026/catalogo_externo_pdg_2026.csv} y su
vista JSON contienen \(2025\) filas: \(1170\) identidades individuales,
\(438\) grupos y \(855\) observables de las tablas sumarias, de los cuales
\(471\) son propiedades de masa y \(384\) de anchura. Once propiedades de
masa se enlazan con reconstrucciones calibradas activas, siete conservan
únicamente una decodificación inversa histórica y las \(2007\) filas
restantes carecen de realización HMT--MD sellada. Ninguna anchura se
promueve a resultado interno. Este conteo heterogéneo describe exclusivamente
la cobertura documental del inventario congelado y no define el dominio de
HMT--MD ni una frontera epistemológica para sus familias internas. El cardinal
documental \(2025\), el año de una edición bibliográfica y la identidad
independiente \(45^2=2025\) de APP son tres datos no relacionados.

\begin{table}[htbp]
\centering
\caption{Campos causales obligatorios del catálogo digital.}
\label{tab:esquema-catalogo-tipado}
\small
\begin{tabularx}{\textwidth}{@{}L{0.25\textwidth}Y@{}}
\toprule
Campo & Valores y función\\
\midrule
\texttt{causal\_direction} &
\texttt{FORWARD}, \texttt{INVERSE} o \texttt{NONE}; indica el sentido de la
aplicación, no su calidad empírica.\\
\texttt{construction\_mode} &
\texttt{TARGET\_BLIND}, \texttt{TARGET\_RESIDUAL\_CALIBRATED} o
\texttt{EXTERNAL\_ONLY}; declara si el blanco intervino en la selección.\\
\texttt{coverage\_state} &
\texttt{ACTIVE}, \texttt{CANDIDATE}, \texttt{ABSENT} o
\texttt{HISTORICAL}; no sustituye al estatuto probatorio.\\
\texttt{statute} &
Uno de los seis estatutos canónicos; la fuente PDG conserva siempre
\texttt{RECONOCIMIENTO\_EXTERNO}.\\
\texttt{observable\_kind} &
Identidad, masa, anchura, diferencia cuadrática, carga topológica,
dimensión cuántica, brecha espectral u otro observable tipado.\\
\texttt{unit}, \texttt{scheme}, \texttt{scale} &
Unidad, esquema y escala de renormalización cuando proceden; se usa
\texttt{NOT\_APPLICABLE} o \texttt{UNAVAILABLE} y nunca un valor supuesto.\\
\texttt{covariance} &
Covarianza publicada; puede ser \texttt{UNAVAILABLE} o
\texttt{NOT\_APPLICABLE}, pero no se reconstruye artificialmente.\\
\bottomrule
\end{tabularx}
\end{table}

\section{A. Atlas HMT interno completo}

\begin{longtable}{@{}llrrp{0.27\textwidth}@{}}
\caption{Los trece tipos celulares del atlas interno.}
\label{tab:atlas-trece-tipos}\\
\toprule
Sector & Tipo celular & Celdas & Total sector & Estatuto y alcance\\
\midrule
\endfirsthead
\toprule
Sector & Tipo celular & Celdas & Total sector & Estatuto y alcance\\
\midrule
\endhead
\(\ell^-\)&\(G0\)&1&25&Centro único \((5,5)\);
\textsc{exacto interno}.\\
\(\ell^-\)&\(G2\)&8&&Capa combinatoria; interfaz física abierta.\\
\(\ell^-\)&\(G4\)&16&&Capa combinatoria; interfaz física abierta.\\
\(\nu\)&\(G1\)&2&20&Capa neutra; no fija una masa de neutrino.\\
\(\nu\)&\(G2\)&4&&Capa neutra; no fija una masa de neutrino.\\
\(\nu\)&\(G3\)&6&&Capa neutra; no fija una masa de neutrino.\\
\(\nu\)&\(G4\)&8&&Capa estéril/exótica candidata.\\
\(d\)-tipo&\(G1\)&2&20&Tipo celular; requiere cociente físico.\\
\(d\)-tipo&\(G2\)&4&&Tipo celular; requiere cociente físico.\\
\(d\)-tipo&\(G3\)&6&&Tipo celular; requiere cociente físico.\\
\(d\)-tipo&\(G4\)&8&&Tipo celular; requiere cociente físico.\\
\(u\)-tipo&\(G1\)&4&16&Tipo celular; requiere cociente físico.\\
\(u\)-tipo&\(G3\)&12&&Tipo celular; requiere cociente físico.\\
\bottomrule
\end{longtable}

El censo satisface
\[
 81=25_{\ell^-}+20_\nu+20_{d\text{-tipo}}+16_{u\text{-tipo}},
 \qquad
 \mathcal C_{81}\longrightarrow\mathcal R_{56}.
\]
El cociente ulterior produce trece multisecciones. El electrón central es la
única excepción puntual; las otras doce clases no admiten un selector
equivariante desde la etiqueta gruesa sola. En el sector quark-tipo, el
residuo \(r-c\bmod3\) reparte exactamente
\[
 36=12+12+12.
\]
Fijadas la carta qutrit, su base, su orientación y el carácter de fase, este
balance alimenta la construcción exacta de
\(\mathfrak{sl}_3(\mathbb C)\), de su forma real compacta y de una teoría gauge
finita en el \cref{chap:algebra-color-qutrit}. La identificación con el color
físico y el límite cromodinámico permanecen abiertos.

\begin{table}[htbp]
\centering
\caption{Ocho alturas publicadas de las torres globales.}
\label{tab:ocho-torres-globales}
\begin{tabular}{@{}ccl@{}}
\toprule
Altura & Descomposición en \(\langle90,120\rangle\) & Función\\
\midrule
90&\(90\)&generador\\
120&\(120\)&generador\\
180&\(2\cdot90\)&refinamiento\\
210&\(90+120\)&refinamiento\\
240&\(2\cdot120\)&refinamiento\\
270&\(3\cdot90\)&refinamiento\\
360&\(3\cdot120\)&refinamiento\\
420&\(2(90+120)\)&refinamiento\\
\bottomrule
\end{tabular}
\end{table}

La completitud de las torres sobre \(\mathbb R_{>0}\) es reconstructiva:
permite aproximar un blanco positivo mediante el retículo logarítmico
declarado, pero no selecciona prospectivamente una especie y no contiene el
valor nulo.

\section{B. Partículas elementales}

\begin{longtable}{@{}L{0.19\textwidth}L{0.23\textwidth}L{0.20\textwidth}L{0.28\textwidth}@{}}
\caption{Vista impresa por familias de partículas elementales.}
\label{tab:catalogo-elementales}\\
\toprule
Familia externa & Observable metrológico & Cobertura HMT--MD &
Fuerza probatoria\\
\midrule
\endfirsthead
\toprule
Familia externa & Observable metrológico & Cobertura HMT--MD &
Fuerza probatoria\\
\midrule
\endhead
Leptones cargados \(e,\mu,\tau\)&
Masa de reposo o polo, carga y anchura cuando procede&
Centro interno para \(e\); masas de \(\mu,\tau\) en tabla calibrada&
Identidades y representaciones: \textsc{interfaz física}; masas tabuladas:
\textsc{reconstrucción calibrada}.\\
Neutrinos \(\nu_e,\nu_\mu,\nu_\tau\)&
Diferencias de masa al cuadrado, ordenamiento y cotas&
Capas neutras internas, sin tres masas inventadas&
\textsc{reconocimiento externo}; selector físico
\textsc{programa abierto}.\\
Quarks \(u,d,s,c,b,t\)&
Masa con esquema y escala; anchura para \(t\)&
Tipos \(u/d\), color residual y decodificación histórica parcial&
No hay flujo de renormalización en \(\chi_{\rm mass}\);
realización \textsc{abierta}.\\
Fotón \(\gamma\)&
Sector físicamente sin masa y dos polarizaciones&
Sección nula candidata; ninguna fila de masa positiva&
\textsc{interfaz física}/\textsc{programa abierto}; el cero externo no es
salida de \(\chi_{\rm mass}>0\).\\
Gluones \(g\)&
Sector gauge sin masa; color adjunto&
Sin representación HMT sellada&
\texttt{ABSENT}; \textsc{programa abierto}.\\
Bosones \(W^\pm,Z\)&
Polo complejo, anchura y convención de masa&
Propiedades de masa en tabla calibrada; anchuras no construidas&
\textsc{reconstrucción calibrada} para la masa;
\textsc{reconocimiento externo} para la anchura.\\
Higgs \(H\)&
Polo, anchura y convención experimental&
Decodificador inverso histórico, sin salida prospectiva&
\texttt{HISTORICAL}; comparación externa.\\
\bottomrule
\end{longtable}

\section{C. Hadrones, resonancias y exóticos compuestos}

\begin{longtable}{@{}L{0.20\textwidth}L{0.22\textwidth}L{0.24\textwidth}L{0.25\textwidth}@{}}
\caption{Vista impresa de objetos compuestos.}
\label{tab:catalogo-compuestos}\\
\toprule
Clase & Observable & Operación HMT--MD requerida & Estado\\
\midrule
\endfirsthead
\toprule
Clase & Observable & Operación HMT--MD requerida & Estado\\
\midrule
\endhead
Mesones estables o longevos&
Masa de polo/reposo y vida media&
Empalmar extensiones y registros antes de extraer la firma&
\(\pi,K\) poseen filas calibradas; el resto es comparación externa.\\
Resonancias mesónicas&
Polo complejo y anchura&
Carácter con codominio complejo o par masa--anchura, aún no construido&
\textsc{interfaz física abierta}.\\
Bariones&
Masa, anchura, carga y composición&
Composición enriquecida anterior a la proyección&
\(p,n\) poseen filas calibradas; no se suman firmas de constituyentes.\\
Tetraquarks y pentaquarks&
Polo, anchura, números cuánticos y canal de decaimiento&
Árbol de composición y multisección enriquecida&
\textsc{reconocimiento externo}; cobertura \texttt{ABSENT}.\\
Glueballs e híbridos&
Espectro y mezcla dependientes del modelo/retículo&
Representación gauge y regla de mezcla&
Candidatos; \textsc{programa abierto}.\\
\bottomrule
\end{longtable}

\section{D. Sectores topológicos}

\begin{longtable}{@{}L{0.18\textwidth}L{0.24\textwidth}L{0.24\textwidth}L{0.25\textwidth}@{}}
\caption{Invariantes exigidos para sectores topológicos.}
\label{tab:catalogo-topologico}\\
\toprule
Sector & Invariantes pertinentes & Resultado interno & Realización\\
\midrule
\endfirsthead
\toprule
Sector & Invariantes pertinentes & Resultado interno & Realización\\
\midrule
\endhead
\(\mathbb Z_6\) abeliano&
Carga de trenza, carácter y fase&
Lectores módulo \(2\) y \(3\) de \(B_n^{\rm ab}\simeq\mathbb Z\)&
Cuatro caracteres no equivalen, sin más, a cuatro fases físicas.\\
Centro apuntado APP&
Carga--flujo, \(qdim\), dimensión total&
\(Z(\mathrm{Vec}_{(\mathbb Z/9)^2})\): \(6561\) simples,
\(qdim=1\), \(\mathcal D=81\)&
Modelo matemático apuntado; no realiza Fibonacci ni Ising.\\
Fibonacci&
\(\tau\), fusión, \(qdim\), matrices \(F/R\), giro, brecha y
desdoblamiento&
Anillo basado asociado a \(F_{\rm av}\), con
\(\tau^2=1+\tau\) y \(\operatorname{FPdim}\tau=\varphi\)&
La categorificación trenzada es candidata; no existe masa universal en
MeV.\\
Ising/Majorana&
\(1,\psi,\sigma\), \(\sigma^2=1+\psi\), \(d_\sigma=\sqrt2\), paridad,
\(F/R\), brecha y desdoblamiento&
No se sigue de \(C_M=-\mathrm{rev}\) ni de CAR&
Modo cero dependiente de plataforma; \textsc{programa abierto}.\\
\bottomrule
\end{longtable}

\section{E. Candidatos BSM}

\begin{longtable}{@{}L{0.20\textwidth}L{0.24\textwidth}L{0.25\textwidth}L{0.22\textwidth}@{}}
\caption{Clases BSM incluidas como inventario abierto, no como predicciones.}
\label{tab:catalogo-bsm}\\
\toprule
Clase & Datos mínimos & Estructura HMT--MD pendiente & Estado\\
\midrule
\endfirsthead
\toprule
Clase & Datos mínimos & Estructura HMT--MD pendiente & Estado\\
\midrule
\endhead
Majorana relativista&
Representación autorconjugada y término de masa&
Entrelazador relativista distinto de \(C_M\)&
\texttt{CANDIDATE}; \textsc{programa abierto}.\\
Neutrinos estériles&
Mezcla, diferencias cuadráticas y cotas&
Selector prospectivo de la cuarta capa neutra&
\texttt{CANDIDATE}; sin masa inventada.\\
Axiones o ALP&
Masa, acoplamientos y escala&
Representación y carácter de acoplamiento&
\texttt{ABSENT}; comparación externa.\\
Monopolos&
Carga magnética, masa y topología&
Sector de haz y cuantización de carga&
\texttt{ABSENT}; interfaz abierta.\\
Taquiones&
Dispersión, estabilidad y causalidad&
Ningún modo físico taquiónico admitido por el criterio gravitatorio&
Control negativo, no especie predicha.\\
\bottomrule
\end{longtable}

\section{Protocolo prospectivo y evaluación fibra a fibra}

Para cada caso se congela la clave
\[
 (\texttt{blind\_case\_id},\texttt{candidate\_id},
   \texttt{observable\_kind})
\]
antes de abrir la magnitud externa. La cadena autorizada es
\[
\begin{aligned}
&\text{descriptores no metrológicos}\\
&\dashrightarrow
\{\text{multisecciones completas compatibles}\}_{\rm sellado}\\
&\longrightarrow
\text{para cada extensión de cada multisección: ruta observable}\\
&\longrightarrow\text{registro}
\longrightarrow\text{firma}
\longrightarrow\chi_{\rm form}
\longrightarrow\text{especialización}
\longrightarrow\text{observable}.
\end{aligned}
\]
La salida del emparejamiento físico sigue siendo la multisección enriquecida;
la evaluación posterior se ejecuta fibra a fibra. Si el selector consulta la
masa, la anchura o una firma residual obtenida del blanco, la salida se
rotula \textsc{reconstrucción calibrada}. Para convertir el protocolo en una
predicción ejecutada se exige, además, un entorno sin acceso a blancos,
sellado previo del multiconjunto, regla de desempate fijada y apertura
posterior de las magnitudes.

\begin{criterion}[Integridad del catálogo tipado]
\label{crit:integridad-catalogo-tipado}
El catálogo queda invalidado si confunde identidad con propiedad de masa,
presenta una decodificación inversa como flecha prospectiva, asigna masa
universal a una carga topológica, suma firmas de constituyentes sin empalmar
antes las extensiones, omite esquema o escala de un quark, o atribuye al
carácter positivo un sector de masa exactamente nula.
\end{criterion}
